%------------------------------------------------------------------------------%
\section{Função Afim ou do 1\textsuperscript{o} grau}

Não sei se você já observou, mas o preço a pagar por uma corrida de táxi
depende do número de quilômetros que se percorre, além disso, adicionamos a
esse valor uma taxa fixa, a chamada ``bandeirada''. Se considerarmos, por
exemplo, que a bandeirada custa \reais{4.70} e cada quilômetro rodado custa
\reais{0.90}, então o valor a ser pago por uma corrida de táxi é dado pela
função
\[
  f(x) = 0,90x + 4,70.
\]
Essa função é um exemplo de
\emph{função afim}\index{Função!afim},
que será estudada neste capítulo.

Uma função é dita \emph{Afim} (ou do 1\textsuperscript{o} grau)
se puder ser escrita na forma
\begin{equation}
  \label{eqretageral}
  f(x)=ax+b,
\end{equation}
em que $a, b \in \R$, $a\neq 0$.Uma função quadrática é um polinômio de grau 2:
\

As funções
\[
  y = 3x-5,     \qquad
  y = 3x        \qquad\text{ e }\qquad
  y = -6x-1989
\]
são exemplos de funções do 1\textsuperscript{o} grau.

Os coeficientes $a$ e $b$ recebem o nome respectivamente de
\emph{coeficiente angular}\index{Coeficiente!angular} e
\emph{coeficiente linear}\index{Coeficiente!linear}.
O coeficiente angular determina a
inclinação da reta e o coeficiente linear determina o lugar em que a reta toca
o eixo $y$ uma vez que $(0, b)\in G(f)$.
Considere uma função $f(x)=ax+b$ e
sejam $P$ e $Q$ dois pontos quaisquer do gráfico de $f$, ou seja,
$P=(x_1, y_1)$ e $Q=(x_2, y_2)$ satisfazendo
$y_1=ax_1+b$ e $y_2=ax_2+b$, conforme o gráfico da Figura~\ref{fig:reta2}.

%------------------------------------------------------------------------------%
\begin{figure}
  \centering
  \includegraphics[width=0.6\textwidth,trim={0 0 9cm 1cm},clip]{reta2}
  \caption{???}
  \caption{Representação gráfica de uma função afim e a
          inclinação da reta ($\theta$).}
  \label{coeficiente}
\end{figure}
%------------------------------------------------------------------------------%

Segue que a divisão da variação em $y$ pela variação em $x$ nos fornece
\begin{equation}\label{eqdelta}
  \frac{\Delta y}{\Delta x}
  = \frac{y_2-y_1}{x_2-x_1}
  = \frac{ax_2-ax_1}{x_2-x_1}
  = a.
\end{equation}
Isso quer dizer que podemos obter o valor do coeficiente angular da reta a
partir de dois pontos quaisquer da reta.
Conforme representado na Figura \ref{coeficiente}, essa razão representa
a tangente do ângulo $\theta$ dado pela inclinação da reta em relação
ao eixo $x$.
Se $(x, y)$ é um ponto genérico dessa reta, então
\[
  \frac{y-y_1}{x-x_1}=a
\]
que pode ser reescrito como
\begin{gather*}
  y-y_1 = a(x-x_1)  \nonumber \\
  y = a(x-x_1)+y_1. \label{eqdareta}
\end{gather*}
Quando $b=0$, dizemos que a função é \emph{Linear}\index{Função!linear}.

%------------------------------------------------------------------------------%
\subsection{Como determinar a equação de uma reta dado dois pontos}

Levando em consideração qualquer reta (diferente das retas $x=c$, com $c$
constante) é determinada pela equação \eqref{eqretageral}, que é a equação da
reta com coeficiente angular $a$ e coeficiente linear $b$. Além disso, dois
pontos determinam uma única reta, então se tivermos dois pontos de uma reta
$P=(x_1, y_1)$ e $Q=(x_2, y_2)$ podemos determinar essa equação de duas maneiras
práticas:

\emph{1\textsuperscript{o} modo:}
Como $P=(x_1, y_1)$ e $Q=(x_2, y_2)$ pertencem a reta,
então suas coordenadas satisfazem a equação \ref{eqretageral}, com isso podemos
encontrar $a$ e $b$ resolvendo o sistema
\[
  \begin{cases}
    ax_1+b=y_1 \\
    ax_2+b=y_2
  \end{cases}
\]

\emph{2\textsuperscript{o} modo:}
Encontramos o coeficiente angular através a
expressão
\[
  a = \frac{y_2-y_1}{x_2-x_1}
\]
e depois encontramos o coeficiente linear $b$ através de uma da
equação $y=ax+b$, substituindo um dos pontos dados
$P$ e $Q$.

%------------------------------------------------------------------------------%
\begin{observation}
  Se $r$ e $s$ são duas retas paralelas então ambas possuem o mesmo
  coeficiente angular.
\end{observation}
%------------------------------------------------------------------------------%

%------------------------------------------------------------------------------%
\begin{observation}
  Se $r$ e $s$ são duas retas perpendiculares, sejam $a_r$ e $a_s$ os
  coeficientes angulares das retas $r$ e $s$ respectivamente. Temos que
  \[
    a_r=-\frac{1}{a_s}.
  \]
\end{observation}
%------------------------------------------------------------------------------%

%------------------------------------------------------------------------------%
\begin{example}
  Determine a equação da reta que passa pelos pontos $P=(2,1)$ e $Q=(3,-1)$.
  (Faça dos dois modos).

  \solution
  \emph{1\textsuperscript{o} modo:}

  Como $P$ e $Q$ pertencem a reta que estamos
  querendo encontrar ($y=ax+b$), então precisamos resolver o sistema
  \[
    \begin{cases}
      2a+b=1 \\
      3a+b=-1
    \end{cases}
  \]
  cuja solução é $a=-2$ e $b=5$, logo a equação da reta pedida é $y=-2x+5$.

  \emph{2\textsuperscript{o} modo:}

  O coeficiente angular da reta pedida é
  \[
    a
    = \frac{y_2-y_1}{x_2-x_1}
    = \frac{-1-1}{3-2}=-2.
  \]
  Como a reta passa pelos pontos $P$ e $Q$, podemos escolher qualquer um desses
  pontos para substituir na equação da reta $y=-2x+b$, para encontrarmos o valor
  de $b$. Substituindo o ponto $P$, temos que
  \[
    1 = -2 \cdot 2+b
  \]
  então $b=5$ e a equação da reta pedida é
  \[
    y=-2x+5.
  \]
\end{example}
%------------------------------------------------------------------------------%

%------------------------------------------------------------------------------%
\begin{example}\label{ex1}
  O número de veículos de uma cidade cresceu linearmente a partir do ano
  2000. Sabendo que a cidade tinha 150 mil veículos em 2004 e 210 mil
  veículos em 2012.

  \solution
  Defina uma função que forneça o número de veículos (em milhares)
  em relação ao tempo (anos) transcorrido desde o ano 2000.

  Inicialmente vamos considerar que $x=0$ corresponde ao ano 2000, $x=1$
  corresponde ao ano 2001, $x=2$ corresponde ao ano 2002 e assim
  sucessivamente. Seja $y$ a quantidade de veículos (cada mil) dessa cidade
  e que cresce linearmente (a função é uma reta). Temos dois pontos dados
  no enunciado $P=(4, 150)$, que quer dizer que em $x=4$ (ano de 2004) a
  cidade tinha 150 mil veículos e $Q=(12, 210)$, que quer dizer que em
  $x=12$ (ano de 2012) a cidade tinha 210 mil veículos. Logo, o
  coeficiente angular dessa reta será
  \[
    a
    = \frac{y_2-y_1}{x_2-x_1}
    = \frac{210-150}{12-4}
    = \frac{60}{8}
    = \frac{15}{2}
    = 7,5
  \]
  Para encontrarmos o valor de $b$, substituímos o ponto $P$
  (poderia ser o $Q$, tanto faz)
  na equação
  \[
    y = 7,5x + b,
  \]
  logo
  \[
    150 =4 \ddot 7,5 + b
  \]
  então $b=120$, logo a equação da reta pedida é
  \[
    y = 7,5x + 120.
  \]

  Determine aproximadamente o número de veículos no ano 2000.

  O ano 2000 corresponde a $x=0$, então
  \[
    y = 7,5 \cdot 0 + 120 = 120.
  \]
  Logo em 2000 o número de veículos era 120 mil.

  Determine em que ano a cidade tera 360 mil veículos.

  Nesse caso, basta substituirmos $y$ por 360, na equação encontrada, logo
  \[
    360 = 7,5x + 120
  \]
  então $x=32$. Logo, a cidade terá 360 veículos no ano 2032.
\end{example}
%------------------------------------------------------------------------------%

%------------------------------------------------------------------------------%
\subsection{Zero e Monotonicidade da função Afim}

Para determinarmos o zero de uma função afim, basta resolver a equação
$ax+b=0$, supondo que $a \neq 0$, temos que o zero de uma função
do 1\textsuperscript{o} grau é
\[
  x = -\frac{b}{a}.
\]

Quanto a monotonicidade de uma função afim se $a>0$ a função é sempre
crescente, se $a<0$ a função é sempre decrescente. Por definição, uma função do
1\textsuperscript{o} grau exige que $a \neq 0$. Caso o coeficiente $a$ seja
igual a zero, a função resulta em $f(x) = b$, sendo classificada apenas como
uma função constante, não possuindo inclinação nem sendo considerada uma função
de 1\textsuperscript{o} grau.

%------------------------------------------------------------------------------%
\begin{figure}
  \centering
  \includegraphics[scale=0.35]{linear1} \qquad
  \includegraphics[scale=0.35]{linear2}
  \caption{Representação gráfica da monotonicidade da função afim,
           destacando o comportamento crescente, $a > 0$, e decrescente,
           $a < 0$, e os respetivos zeros da função \(x = -\frac{b}{a}\).}
\end{figure}
%------------------------------------------------------------------------------%

%------------------------------------------------------------------------------%
\subsection{Gráfico de uma função Afim}

Como já sabemos que o gráfico de uma função do 1\textsuperscript{o} grau
é uma reta,
precisamos de apenas dois pontos pertencentes a essa reta para o esboço do
gráfico. Em particular, se tivermos os interceptos com os eixos coordenados
\[
  \left( -\frac{b}{a}, 0 \right)
  \quad\text{ e }\quad
  (0, b)
\]
conseguimos fazer o esboço do gráfico com muita facilidade.

%------------------------------------------------------------------------------%
\begin{example}
  Construa o gráfico da função $f(x)=2x-6$.

  \solution
  \emph{Crescimento}

  Como $a = 2$ ($a > 0$), a função é crescente.

  \emph{Zero da Função (Intercepto no eixo $x$)}

  Fazemos $f(x) = 0$
  \begin{align*}
    2x - 6 & = 0 \\
    2x     & = 6 \\
    x      & = 3
  \end{align*}
  O intercepto no eixo das abscissas é o ponto $(3, 0)$.

  \emph{Intercepto no eixo $y$}

  É o valor do coeficiente linear $b$. Como $b = -6$, o ponto é $(0, -6)$.

  \IncludeGraphicCentredEx{rlin1}[0.5]
\end{example}
%------------------------------------------------------------------------------%

%------------------------------------------------------------------------------%
\begin{example}
  Construa o gráfico da função $g(x)=-x-4$.

  \solution
  \emph{Crescimento}

  Como $a = -1$ ($a < 0$), a função é decrescente.

  \emph{Zero da Função (Intercepto no eixo $x$)}

  Fazemos $g(x) = 0$:
  \begin{align*}
    -x - 4 & =  0 \\
    -x     & =  4 \\
    x      & = -4
  \end{align*}
  O intercepto no eixo das abscissas é o ponto $(-4, 0)$.

  \emph{Intercepto no eixo $y$}

  Como $b = -4$, o ponto é $(0, -4)$.

  \IncludeGraphicCentredEx{rlin2}[0.5]
\end{example}
%------------------------------------------------------------------------------%

%------------------------------------------------------------------------------%
\begin{example}
  Construa o gráfico da função $h(x)=2x$.

  \solution
  \emph{Crescimento}

  Como $a = 2$ ($a > 0$), a função é crescente.

  \emph{Zero da Função (Intercepto no eixo $x$)}

  \begin{align*}
    2x & = 0 \\
     x & = 0.
  \end{align*}
  O ponto é $(0, 0)$.

  \emph{Intercepto no eixo $y$}

  Como $b = 0$, o ponto também é $(0, 0)$.

  \emph{Nota Didática}

  Como a reta passa pela origem, utiliza-se um ponto auxiliar,
  como $x=1$, $h(1)=2$, resultando no ponto $(1, 2)$.

  \IncludeGraphicCentredEx{rlin3}[0.5]
\end{example}
%------------------------------------------------------------------------------%

%------------------------------------------------------------------------------%
\subsection{Algumas retas particulares:}

%------------------------------------------------------------------------------%
\begin{figure}
  \begin{center}
  \includegraphics[scale=0.3]{cons1} \qquad
  \includegraphics[scale=0.3]{cons2}
  \end{center}
  \caption{Gráfico da reta horizontal $y=b$ e da reta vertical $x=a$.}
\end{figure}
%------------------------------------------------------------------------------%
